% swift-equatable-hashable-comparable.tex — the equality/hash contract,
% synthesis rules, hash(into:) + Hasher, mutable reference keys, Comparable,
% Identifiable vs Hashable, NSObject isEqual/hash bridging.
% Sources (checked 2026-10-06):
%   swift-evolution SE-0185 (Equatable / Hashable synthesis, 4.1) · SE-0206 (hash(into:),
%     Hasher = "SipHash-1-3", per-run seed, SWIFT_DETERMINISTIC_HASHING, 4.2) · SE-0261
%     (Identifiable, 5.1) · SE-0266 (enum Comparable incl. payloads that are all Comparable;
%     not raw-value enums, 5.3) · SE-0015 (tuple comparison operators, 2.2 — stdlib == / <
%     for tuples up to arity 6; tuples themselves do NOT conform: SE-0283 was returned for
%     revision) — checked against the evolution index and the 0206 text
%   developer.apple.com/documentation/swift/hashable · /objectivec/nsobject/isequal(_:) ·
%     /swift/identifiable (AnyObject default id = ObjectIdentifier) · /swift/range/clamped(to:)
% @labels: area=swift kind=concept platform=apple topic=language,data discussion=156
% @tags: equatable, hashable, comparable, hasher, hash-into, siphash, hash-contract, identifiable, objectidentifier, isequal, nsobject
\documentclass[8pt]{extarticle}
\usepackage{printup-sheet}
\usepackage{array}

\lstdefinelanguage{SwiftSheet}{
  morekeywords={protocol,class,final,struct,enum,func,var,let,init,case,static,
    override,inout,extension,
    if,else,return,guard,self,nil,private,some,switch,default,in,true,false,Int,String,Bool},
  sensitive=true, morecomment=[l]{//}, morecomment=[s]{/*}{*/}, morestring=[b]",
  literate={??}{{\hbox{?}\hbox{?}}}2 {==}{{\hbox{=}\hbox{=}}}2
           {!=}{{\hbox{!}\hbox{=}}}2 {->}{{\hbox{-}\hbox{>}}}2 {===}{{\hbox{=}\hbox{=}\hbox{=}}}3}

% 0xProto ligatures: one \hbox per character, for use INSIDE \texttt
\newcommand\EQ{\hbox{=}\hbox{=}}
\newcommand\EQQ{\hbox{=}\hbox{=}\hbox{=}}
\newcommand\NE{\hbox{!}\hbox{=}}
\newcommand\LE{\hbox{<}\hbox{=}}
\newcommand\GE{\hbox{>}\hbox{=}}
\newcommand\arr{\hbox{-}\hbox{>}}
\newcolumntype{L}[1]{>{\raggedright\arraybackslash}p{#1}}

\tikzset{
  bk/.style={draw=sheetGrey, minimum width=6.5mm, minimum height=5mm, font=\tiny, inner sep=0},
  el/.style={box, font=\ttfamily\tiny, inner sep=1.2pt, minimum height=3.8mm},
  bad/.style={el, draw=sheetRed, fill=sheetRed!7},
  good/.style={el, draw=sheetGreen, fill=sheetGreen!10},
  lbl/.style={font=\tiny, text=black!75, inner sep=1pt, align=center},
  hd/.style={font=\bfseries\small, anchor=west},
  hs/.style={sb, font=\ttfamily\scriptsize},
  sb/.style={box, font=\scriptsize, inner sep=1.5pt, minimum height=5mm},
}

\begin{document}

\sheettitle{Equatable · Hashable · Comparable — the contracts}{swift · folio}

\oneliner{\textbf{Equatable} = value equality (\texttt{\EQ}); \textbf{Hashable} adds
\texttt{hash(into:)} under one rule — \textbf{\texttt{a \EQ{} b} $\Rightarrow$ equal hashes}
(never the converse); \textbf{Comparable} adds a strict total order from \textbf{only
\texttt{<}}. \texttt{Set}/\texttt{Dictionary} find a key by \emph{hash first, then \texttt{\EQ}}
— break the contract or mutate a key and elements vanish.}

\vspace{3pt}
\noindent\begin{tikzpicture}[sheet]
  \foreach \x in {4.6,9.7} \draw[sheetGrey!40] (\x,1.45) -- (\x,-1.95);
  % ── A: Hasher pipeline ──
  \node[hd] at (0,1.3) {\textcolor{sheetBlue}{A} how a lookup works};
  \node[hs] (k) at (0.55,0.65) {key};
  \node[hs, text width=17mm] (h) at (2.3,0.65) {Hasher(seed)\\combine(id)\\finalize()};
  \node[hs] (i) at (3.95,0.65) {Int};
  \draw[flow] (k) -- (h); \draw[flow] (h) -- (i);
  \foreach \b in {0,...,5} \node[bk] (b\b) at (0.6+\b*0.68,-0.55) {\b};
  \draw[hot] (i.south) -- node[lbl, right]{\texttt{\% count}} (b3.north);
  \node[lbl, text=sheetBlue] at (2.3,-1.05) {probe bucket 3, then \texttt{\EQ} on each candidate};
  \node[lbl, text=sheetBrown] at (2.3,-1.55) {seed = random per \textbf{process}: hashes + Set order\\change every launch (SipHash-1-3)};
  % ── B: broken contract ──
  \node[hd] at (4.7,1.3) {\textcolor{sheetBlue}{B} \texttt{hash} uses MORE than \texttt{\EQ}};
  \node[lbl, anchor=west] at (4.7,0.85) {\texttt{\EQ} compares \texttt{id}; \texttt{hash} combines \texttt{id + name}};
  \foreach \b in {0,...,5} \node[bk] (c\b) at (5.2+\b*0.68,0.15) {\b};
  \node[bad, anchor=north] at (c1.south) {7,"Ann"};
  \node[bad, anchor=north] at (c4.south) {7,"Anna"};
  \node[lbl, text=sheetRed] at (7.15,-1.1) {\texttt{a \EQ{} b} but different buckets:\\Set keeps BOTH, lookups miss};
  \node[lbl, text=sheetGreen!60!black] at (7.15,-1.7) {rule: hash a \textbf{subset} of what \texttt{\EQ} reads};
  % ── C: mutated key ──
  \node[hd] at (9.8,1.3) {\textcolor{sheetBlue}{C} mutating a class key in a Set};
  \foreach \b in {0,...,5} \node[bk] (d\b) at (10.3+\b*0.68,0.15) {\b};
  \node[good, anchor=north] (e) at (d2.south) {obj};
  \node[lbl, anchor=west] at (9.8,0.85) {inserted when \texttt{obj.name = "a"} $\to$ bucket 2};
  \draw[hot, draw=sheetRed, dashed] (d5.south) ++(0,-0.1) -- ++(0,-0.35) node[below, lbl, text=sheetRed]{lookup now\\probes 5};
  \node[lbl, anchor=west, text=sheetRed] at (9.8,-1.1) {\texttt{obj.name = "b"} (shared reference!):};
  \node[lbl, anchor=west, text=sheetRed] at (9.8,-1.45) {\texttt{set.contains(obj)} $\to$ \textbf{false}, dups on insert};
  \node[lbl, anchor=west] at (9.8,-1.8) {structs are safe: the Set holds its own copy};
\end{tikzpicture}

\begin{multicols}{2}
\raggedright

\section{How it works}
\begin{itemize}
  \item \textbf{Equatable}: reflexive, symmetric, transitive; \texttt{\NE} is free.
        \textbf{Synthesized} (SE-0185, Swift 4.1) when you declare it and every stored
        property / associated value conforms; payload-less enums get it without
        declaring. Never for \textbf{classes}. Must be declared in the type's own file.
  \item \textbf{Hashable: Equatable}. Implement \texttt{hash(into: inout Hasher)} and
        \texttt{combine} the same fields \texttt{\EQ} reads (SE-0206, Swift 4.2);
        implementing \texttt{hashValue} is deprecated. \texttt{Hasher} = SipHash-1-3,
        randomly seeded per process (HashDoS resistance);
        \texttt{SWIFT\_DETERMINISTIC\_HASHING=1} fixes the seed for tests.
        \textbf{Never persist} a hash value.
  \item \textbf{Comparable: Equatable}. Only \texttt{<} is required; \texttt{>}, \texttt{\LE},
        \texttt{\GE}, \texttt{min}/\texttt{max}, \texttt{sort()}, ranges \texttt{a...b}
        come free. Must be a strict total order consistent with \texttt{\EQ}
        (\texttt{!(a<b) \&\& !(b<a)} $\Rightarrow$ \texttt{a \EQ{} b}).
        Synthesized only for \textbf{enums} without raw values whose payloads are
        Comparable — declaration order (SE-0266, 5.3). Tuples (arity $\le$ 6) compare
        lexicographically: delegate to them.
  \item \textbf{Identifiable} (SE-0261, 5.1): \texttt{var id: ID} where \texttt{ID: Hashable};
        classes get \texttt{ObjectIdentifier(self)} by default. \textbf{Identity} (which row
        is this?) $\neq$ \textbf{equality} (did its content change?). SwiftUI keys state,
        animation and diffing on \texttt{id}.
  \item \textbf{NSObject} bridges: Swift \texttt{\EQ} calls \texttt{isEqual(\_:)},
        \texttt{hashValue} returns \texttt{hash}. Default \texttt{isEqual} is identity.
        Subclass: override \textbf{\texttt{isEqual(\_:)} + \texttt{hash}} — a custom
        \texttt{static \EQ} or \texttt{hash(into:)} is \emph{not} what \texttt{Set},
        \texttt{NSSet} or Foundation call.
\end{itemize}

\section{Example}
\begin{lstlisting}[language=SwiftSheet]
struct User: Hashable {           // identity = id
  let id: Int; var name: String
  static func == (l: User, r: User) -> Bool { l.id == r.id }
  func hash(into h: inout Hasher) { h.combine(id) } // same field
}
struct Version: Comparable {
  let major, minor, patch: Int
  static func < (l: Self, r: Self) -> Bool {       // only <
    (l.major, l.minor, l.patch) < (r.major, r.minor, r.patch)
  }
}
final class Node: Hashable {                        // identity
  static func == (l: Node, r: Node) -> Bool { l === r }
  func hash(into h: inout Hasher) {
    h.combine(ObjectIdentifier(self)) }   // stable
}
\end{lstlisting}

\columnbreak

\section{Which equality for which type?}
{\footnotesize
\begin{tabular}{@{}L{22mm}L{52mm}@{}}
\toprule
\textbf{Type} & \textbf{Choose} \\
\midrule
value struct & synthesize; hand-write only to exclude fields (cache, timestamps) \\
entity (\texttt{User}) & \texttt{\EQ}/hash on \texttt{id}; \texttt{Identifiable} for SwiftUI \\
class, identity & \texttt{\EQQ} + \texttt{ObjectIdentifier} (safe to mutate) \\
class, value-based & only if the hashed fields are \texttt{let} \\
\texttt{NSObject} subclass & \texttt{isEqual(\_:)} + \texttt{hash}, together \\
normalised keys & \texttt{\EQ} and \texttt{hash} on the SAME normal form (lowercased) \\
\bottomrule
\end{tabular}}

\section{Traps}
\begin{itemize}
  \trap{\texttt{hash} combining a field \texttt{\EQ} ignores (panel B) — equal values,
        different buckets.}
  \trap{Mutable class with value-based hash inside a \texttt{Set}/as a dict key (panel C).}
  \trap{\texttt{isEqual} overridden, \texttt{hash} forgotten: default hash is identity
        $\to$ broken \texttt{NSSet}/\texttt{Set}.}
  \trap{\texttt{Set}/\texttt{Dictionary} order differs per launch — never assert or
        persist it.}
  \trap{\texttt{Double.nan \EQ{} .nan} is \texttt{false}: breaks reflexivity,
        \texttt{contains(.nan)} fails, \texttt{sort} with NaN is unspecified.}
  \trap{\texttt{ForEach(items, id: \textbackslash.self)} = identity is the whole value:
        editing a row \emph{changes its identity} (state reset, wrong animation);
        duplicates collide. Array index as id breaks on insert/reorder.}
  \trap{Hash collisions are legal: unequal values may share a hash.}
\end{itemize}

\section{Remember}
\textbf{``Equal $\Rightarrow$ same hash; hash $\subseteq$ equality; never mutate a key;
\texttt{<} buys all of Comparable; \texttt{id} is who, \texttt{\EQ} is what.''}


\end{multicols}

\noindent{\footnotesize\color{sheetGrey}\textit{Related:} value-vs-reference ·
swift-enums-pattern-matching (synthesis) · strings-and-collections (hash tables) ·
swiftui-performance-identity · tableview-diffable (item identifiers)}

\end{document}
